\sectionright{\sffamily Signals, Routing and Signal Manipulation}

%%%%%%%%%%%%%%%%%%%%%%%%%%%%%%%%%%%%%%%%%%%%%%%%%%%%%%%%%%%%%%%%%%%%%%%%%%%%%%%%%%%%%%%%%%%%%%%%%%%%%%%%%%%%%%%%%%%%%%%%%%%%%%%%%%%%%%%%%%%%%%%%%%%%%%%%%%%%%%%%%%%%%%%%%
\subsection{\sffamily Signal Manipulation}

    \subsubsection{\sffamily Description}
        The \texttt{rtl\_sigmanip} element is used for non-arithmetic operations on signals, such as bit-slicing, concatenation, or sign extension. Its elliptical shape and soft orange tint visually separate it from functional logic blocks, indicating that the data is being reinterpreted or rearranged rather than computed.

    \subsubsection{\sffamily Attributes}
    \begin{center}
        \begin{tabular}{lp{8cm}}
        \toprule
            \textbf{Attribute} & \textbf{Effect} \\
            \midrule
            \texttt{inner sep} & Adjusts the padding around the text (Default: 2pt). \\
            \bottomrule
        \end{tabular}
    \end{center}

    \subsubsection{\sffamily Listing}
\begin{RTLexample}[Signal Slicing]
\begin{tikzpicture}[>=Stealth, thick, scale=0.8]
    \node[rtl_reg] (R1) at (0,0) {Reg};
    \node[rtl_sigmanip, right=1.5cm of R1] (Slice) {[15:8]};
    
    \draw[->, rtl_bus={16}] (R1) -- (Slice);
    \draw[->, rtl_bus={8}] (Slice) -- ++(3.5,0) node[right] {Byte\_1};
\end{tikzpicture}
\end{RTLexample}

%%%%%%%%%%%%%%%%%%%%%%%%%%%%%%%%%%%%%%%%%%%%%%%%%%%%%%%%%%%%%%%%%%%%%%%%%%%%%%%%%%%%%%%%%%%%%%%%%%%%%%%%%%%%%%%%%%%%%%%%%%%%%%%%%%%%%%%%%%%%%%%%%%%%%%%%%%%%%%%%%%%%%%%%%
\subsection{\sffamily Bus Styling}

    \subsubsection{\sffamily Description}
        To represent multi-bit data paths, \texttt{register-transfer-level} provides the \texttt{rtl\_bus} decoration. Applying this style to a path increases the line width and adds a diagonal slash at the midpoint, accompanied by a label indicating the bus width. 
        
        A key feature of this decoration is its geometric intelligence: the element automatically calculates the total path length and places the \texttt{/} marker at exactly half the distance (\texttt{pos=0.5}). This ensures that the bus indicator remains perfectly centered, regardless of whether the path is a simple straight line or a complex route with multiple segments (e.g., a Z-route).
        Alternatively the position of the bus decoration can be set manually by using the \texttt{rtl\_bus\_flex} element.
    \subsubsection{\sffamily Attributes}
    \begin{center}
        \begin{tabular}{lp{8cm}}
        \toprule
            \textbf{Attribute} & \textbf{Effect} \\
            \midrule
            \texttt{rtl\_bus=\{width\}} & Activates the bus decoration and sets the bit-width label. \\
            \midrule
            \texttt{rtl\_bus\_flex=\{width,pos\}} & Activates the bus decoration, sets the bit-width label and alters the default positions of the decoration. \\
            \bottomrule
        \end{tabular}
    \end{center}

    \subsubsection{\sffamily Listing}
\begin{RTLexample}[Bus Decoration]
\begin{tikzpicture}[>=Stealth, thick]
    \draw[->, rtl_bus={32}] (0,0) -- (3,0) node[right] {Data\_Bus};
    \draw[->, rtl_bus={}] (0,-1) -- (3,-1) node[right] {Generic\_Bus};
    \draw[->, rtl_bus_flex={n,0.33}] (7,0) -- (10,0) node[right] {Flex\_Bus};
\end{tikzpicture}
\end{RTLexample}

%%%%%%%%%%%%%%%%%%%%%%%%%%%%%%%%%%%%%%%%%%%%%%%%%%%%%%%%%%%%%%%%%%%%%%%%%%%%%%%%%%%%%%%%%%%%%%%%%%%%%%%%%%%%%%%%%%%%%%%%%%%%%%%%%%%%%%%%%%%%%%%%%%%%%%%%%%%%%%%%%%%%%%%%%
\subsection{\sffamily Advanced Routing (Z and U Routes)}

    \subsubsection{\sffamily Description}
        Manual routing with standard TikZ coordinates can become tedious in complex diagrams. \texttt{register-transfer-level} introduces two specialized commands to automate common routing patterns:
        
        \begin{itemize}
            \item \textbf{\textbackslash routeZ}: Creates a "Dogleg" or Z-shaped route. It moves horizontally, then vertically, and finally horizontal again to reach the destination.
            \item \textbf{\textbackslash routeU}: Creates a "Feedback" or U-shaped route. This is ideal for signals that need to loop back from an output to a previous input stage.
        \end{itemize}

    \subsubsection{\sffamily Parameters}
    \begin{center}
        \begin{tabular}{lp{8cm}}
        \toprule
            \textbf{Command} & \textbf{Parameters} \\
            \midrule
            \texttt{\textbackslash routeZ} & \texttt{[Style] \{Start\} \{Kink\_X\} \{End\}} \\
            \texttt{\textbackslash routeU} & \texttt{[Style] \{Start\} \{X\_Space\} \{Y\_Height\} \{End\}} \\
            \bottomrule
        \end{tabular}
    \end{center}

    \subsubsection{\sffamily Listing}
\begin{RTLexample}[Advanced Routing]
\begin{tikzpicture}[>=Stealth, thick, scale=0.8]
    \node[rtl_reg] (A) at (0,0) {A};
    \node[rtl_reg] (B) at (4,-2) {B};
    
    % Z-Route with 1.5cm horizontal segment
    \routeZ[->, blue]{A.east}{1.5}{B.west};
    
    % U-Route (Feedback) from B back to A
    \routeU[->, red, rtl_bus={8}]{B.east}{0.8}{4.5}{A.west};
    
    \node[blue, font=\tiny] at (1.5, -0.2) {Z-Route};
    \node[red, font=\tiny] at (2, 2.7) {Feedback (U-Route)};
\end{tikzpicture}
\end{RTLexample}
